\unnumberedchapter{Введение}

Сжатие данных без потерь остаётся одной из самых частых вычислительных
задач: им пользуются архиваторы, сетевые протоколы, файловые системы и
форматы хранения. Скорость архиватора напрямую определяет, сколько
времени пользователь ждёт результата и сколько процессорного времени
расходует сервер. Поэтому даже программа, которая уже работает правильно,
нуждается в оптимизации: тот же результат можно получить в несколько раз
быстрее, если устранить лишнюю работу и использовать возможности
современных процессоров.

Объект работы "--- архиватор MZ77, разработанный в~весеннем семестре
2026~года как курсовая работа по дисциплине «Структуры данных». Программа
написана на языке~C, сжимает и распаковывает файлы алгоритмом
LZ77 и имеет графический интерфейс на GTK~4.
Исходный код хранится в репозитории
\url{https://github.com/mb6ockatf/MZ77}.

Цель работы "--- ускорить сжатие и распаковку в архиваторе MZ77, не изменяя
формат архива: каждый архив, созданный оптимизированной программой, должен
совпадать байт в байт с архивом исходной программы.

Для достижения цели решены следующие задачи:
\begin{enumerate}
\item выполнено профилирование исходной программы, определены её узкие
места и основные метрики; профилирование повторялось после каждого шага
оптимизации;
\item оценена алгоритмическая сложность исходной программы;
\item предложена и реализована алгоритмическая оптимизация: отказ от сдвига
упреждающего буфера после каждого токена;
\item предложены и реализованы четыре машинно-независимые оптимизации:
блочный ввод-вывод, встраивание побитового ввода-вывода, сравнение по
8~байт и оптимизация при компоновке вместе с оптимизацией по профилю;
\item реализованы машинно-зависимые оптимизации: сравнение совпадений
командами AVX-512BW и AVX2 с выбором реализации во время выполнения;
\item оценена алгоритмическая сложность оптимизированной программы;
\item результаты каждого шага измерены на двух компьютерах с процессорами
разных производителей и представлены в виде таблиц и графиков.
\end{enumerate}

Все измерения, их исходные журналы и сценарии для их повторения сохранены в
репозитории проекта (каталог \texttt{bench/log}); порядок воспроизведения
приведён в приложении~А.
